\section{System Design}
\label{sec:design}

\subsection{Architecture and Execution Cycle}
\label{sec:design-overview}

\systemname coordinates session phases, eviction budgets, and restoration timing to reuse GPU space between updates.
It operates above an inference backend, which remains responsible for execution scheduling and model computation.
Batch formation, when used, remains a backend policy; the architecture does not require continuous batching.
A residency planner constructs the plan at admission and revises it as resource demands change.
A session manager follows the plan, using ticks, execution events, and KV readiness to submit work to the backend.
A KV memory manager tracks allocation, valid copies, and references, while a KV transfer engine executes asynchronous copies.
The session manager issues eviction and restoration requests to the memory manager, which submits transfer requests with valid sources and allocated destinations.
A transfer completion reports a finished copy; KV readiness is published only after all required ranges are valid for the corresponding session.
These components separate planning decisions from the safe execution of each cycle.

Consider one session in Figure~\ref{fig:motivation}.
At its tick, historical KV is ready on the GPU; the backend processes new input and generates the next output.
Newly retained KV is backed up to host memory.
After the selected ranges are no longer in use and their host copies are valid, \systemname releases their GPU storage.
Other sessions can use this space until the manager allocates a destination and restores the evicted state before the next tick.
Figure~\ref{fig:comparison}(b) shows this cycle across staggered sessions.

\subsection{Session Groups and Phase Assignment}
\label{sec:design-offsets}

\systemname divides a common period $T$ into $S$ phase slots, with release offsets $\phi_j=jT/S$ for $j=0,\ldots,S-1$.
Sessions assigned to the same slot form a session group and become eligible for execution at the same phase.
Each occupied slot can contain multiple sessions; the design does not require one slot per session.
Sessions retain separate histories and eviction budgets, and the backend forms execution batches from the eligible work.
A session group is therefore a scheduling unit, with no requirement that all members execute as one fixed batch or transfer atomically.

\paragraph{Trading batching efficiency for memory reuse.}
Aligning every session to the same phase makes more work available for batching, but also concentrates demand for resident KV and restoration.
Staggering groups distributes that demand and lets their eviction intervals supply space for one another.
It can reduce batch sizes and increase total execution cost through additional model iterations and weight accesses.
Grouping multiple sessions at each phase preserves opportunities for batching while avoiding a single synchronized burst.
For roughly balanced groups, increasing $S$ reduces the number of sessions available per group from $N$ toward $N/S$.
The planner must choose this tradeoff using the backend's measured execution costs.

The goal is to preserve the service target while reducing peak memory use.
Let $C_{\mathrm{aligned}}$ and $C_{\mathrm{grouped}}$ denote the complete periodic execution costs for the two release patterns at the same history length and workload.
The extra cost $C_{\mathrm{grouped}}-C_{\mathrm{aligned}}$, together with transfer interference and scheduling delays, must fit within the original margin $T-C_{\mathrm{aligned}}$.
The short generation bursts described in Section~\ref{sec:request-classes} explain why a moderate batching penalty can leave useful slack.
Admission checks this margin and each session's completion time under the proposed grouping.
Initial phase alignment and any resulting input delay count toward the service cost.

\subsection{Joint Residency Planning}
\label{sec:design-planning}

The planner chooses phase assignments, per-session eviction amounts, and restoration times over a finite lookahead horizon $H$.
It uses the configured period, retained state, expected growth, execution costs under the proposed groups, host capacity, and effective transfer service.
A free phase slot is insufficient for admission: the proposed sessions must fit the compute budget, individual transfer windows, and instantaneous GPU capacity.
Equation~\ref{eq:memory-gain} shows why transfer utilization alone does not establish a capacity benefit; the planner must reduce peak allocation.

\paragraph{Eviction budgets.}
An eviction budget $e_i$ limits the amount of safely recoverable KV released from session $i$.
It need not cover the entire history.
As Section~\ref{sec:opportunity} shows, full eviction eventually exceeds the transfer budget as retained state grows.
Partial eviction keeps the excess resident and gives the planner control over each restoration's duration and destination size.
A useful budget must free enough memory to matter while leaving a feasible return path.

\paragraph{Restoration timing and space.}
Let $r_{i,k+1}$ be the next tick and $u_i$ the planned restoration start.
Given the effective bandwidth assigned to the copy, the estimate $u_i+e_i/\widehat B_i+\delta_i\le r_{i,k+1}$ must include queueing, completion, and safety allowances in $\delta_i$.
The individual estimates must agree with a feasible shared transfer schedule and Equation~\ref{eq:transfer-budget}.
At every point, total physical GPU allocation must remain within $G_{\mathrm{KV}}$.
This accounting includes the entire destination of each copy from the time it is allocated, with shared blocks counted once.
Starting early can shorten the interval of reclaimed space even when it makes the transfer easier to finish.

The planner therefore checks when earlier groups release space as well as when later groups need it.
It must account for expected history growth throughout $H$; a plan that fits the current state may cease to fit before the horizon ends.
\placeholder{Joint planning algorithm: selection of group count and horizon, session assignment, safe eviction ranges, resource estimates and margins, and the time-and-space feasibility check.}

\subsection{Safe Eviction and Scheduled Restoration}
\label{sec:design-idle}
\label{sec:design-correctness}

\paragraph{Backing and releasing state.}
The memory manager incrementally creates host copies of newly retained KV; valid copies of unchanged history can be reused across periods.
Backing up state does not release its GPU allocation.
Eviction requires both a complete host copy and the absence of conflicting execution or transfer references.
For a shared physical block, all references must permit reclamation.
These conditions preserve the application's selected history while changing where its KV resides.

\paragraph{Allocating and copying.}
\label{sec:design-prefetch}
The session manager requests restoration at the planned time.
The memory manager allocates the full destination before submitting a copy to the transfer engine.
Space reclaimed from another group's idle state can satisfy this allocation.
An ongoing copy already occupies GPU capacity, but its contents become usable only after completion is published.
\placeholder{Choose and justify the restoration allocation boundary: whole-group reservation or per-session allocation under the group schedule; specify readiness and recovery when a later allocation cannot proceed.}

\paragraph{Resuming computation.}
The backend may access restored KV only after all required ranges are valid for the correct session and logical positions.
Concurrent demand, prefetching, and eviction must share a consistent view of allocation and content validity.
If space or transfer service is unavailable at the planned time, the manager must reassess completion before allowing dependent computation.
Demand restoration can recover missing state, but its delay remains part of service latency and can cause a deadline miss.
No timeout permits reading unfinished KV or reusing a referenced block.

\subsection{Plan Reuse and Revision}
\label{sec:design-revisions}

The planner reuses its plan across updates when state growth remains sufficiently predictable over the horizon.
Per-cycle checks enforce current space, reference, and readiness conditions without reselecting every phase and budget.
Accumulated history growth requires infrequent revisions; session arrivals and departures can require larger adjustments.
Keeping some KV resident limits transfer traffic, but that resident portion can still grow and eventually exhaust the available capacity.

Revisions must account for allocations and transfers already in progress.
A new plan cannot release a referenced block or turn an unfinished copy into valid state.
Feasibility within $H$ therefore supports only that planning horizon, with continued service requiring reassessment.
\placeholder{Revision triggers, horizon extension, transition between plans with in-flight work, and service policy when no feasible revision exists.}
